\begingroup
\setlength{\abovedisplayskip}{7pt}
\setlength{\belowdisplayskip}{7pt}
\setlength{\abovedisplayshortskip}{3pt}
\setlength{\belowdisplayshortskip}{4pt}
\def\ixrad{\mathord{\vcenter{\hbox{\Large\(\surd\)}}}}
\def\ixangle#1#2#3{%
  \makebox[1.95em][r]{\(#1\)}%
  \makebox[1.45em][r]{\(#2\)}%
  \makebox[3.10em][r]{\(#3\)}}
\def\ixobsCount#1#2{\makebox[.8em][r]{#1}\enspace\makebox[2.6em][c]{#2}}
\def\ixobs#1#2#3#4#5#6#7#8#9{%
  #1 & \(\left\{\begin{array}{@{}l@{\enspace}l@{}}
  \text{Direct}&\ixangle{#2}{#3}{#4}\\
  \text{Reflected}&\ixangle{#5}{#6}{#7}
  \end{array}\right.\) &
  \begin{tabular}{@{}l@{}}#8\\#9\end{tabular}}

% Printed 32; principal PDF 39; Google PDF 43.
\par\clearpage
\noindent\makebox[12mm][l]{\fontsize{13}{16}\selectfont 32}\parbox[t]{\dimexpr\linewidth-24mm\relax}{\centering\scriptsize\bfseries DETERMINATION OF THE DIFFERENCE IN LATITUDE BETWEEN GÖTTINGEN AND ALTONA.}\makebox[12mm][r]{}\par\medskip

{\small
\setlength{\arraycolsep}{.85em}
\[
\begin{array}{r||r|r||r|r}
\smash{\raisebox{-.7ex}{Star}}&
\multicolumn{2}{c||}{\text{\normalsize\textls[100]{Göttingen.}}}&
\multicolumn{2}{c}{\text{\normalsize\textls[100]{Altona.}}}\\
&\text{Coll. err.}&\text{Weight}&\text{Coll. err.}&\text{Weight}\\
\hline
24&3''{,}93&10{,}91&1''{,}21&12{,}00\\
25&3{,}88&12{,}00&0{,}83&12{,}00\\
26&3{,}54&12{,}00&0{,}76&12{,}00\\
27&4{,}44&4{,}00&1{,}97&7{,}50\\
28&3{,}27&10{,}91&0{,}95&12{,}00\\
29&3{,}89&6{,}00&1{,}15&12{,}00\\
30&4{,}54&10{,}91&1{,}70&12{,}00\\
31&4{,}52&10{,}91&1{,}73&12{,}00\\
32&3{,}26&12{,}00&1{,}13&12{,}00\\
33&3{,}49&10{,}91&1{,}11&10{,}91\\
34&4{,}58&12{,}00&1{,}66&12{,}00\\
35&3{,}61&10{,}00&1{,}37&10{,}00\\
36&4{,}33&12{,}00&1{,}34&10{,}00\\
37&3{,}97&12{,}00&1{,}91&10{,}91\\
38&3{,}52&12{,}00&1{,}56&10{,}00\\
39&4{,}16&3{,}00&1{,}45&8{,}00\\
40&5{,}04&4{,}80&1{,}02&6{,}00\\
41&3{,}46&4{,}80&1{,}75&10{,}00\\
42&3{,}84&4{,}00&2{,}45&8{,}89\\
43&4{,}40&4{,}00&1{,}01&10{,}00
\end{array}
\]
\par
}

The mean values are as follows:
\begin{center}
\begin{tabular}{@{}l@{\qquad}r@{\enspace}l@{}}
Collimation error in Göttingen & \(3''{,}75\) & with weight \(455{,}17\)\\
Collimation error in Altona & \(1{,}40\) & with weight \(432{,}18\).
\end{tabular}
\end{center}

That the collimation error really changed is evident, and there can be
no doubt that the change occurred during transport, even though every
possible precaution was taken.

\begin{center}\textbf{4.}\end{center}

Although the result obtained for the difference in latitude is entirely
satisfactory, it is worth observing, at least on theoretical grounds,
that the combination of the observations used in Art. 1 is still not
the most advantageous possible unless every star at each site was observed
as often in one position of the sector as in the other. In fact, the
determination of the true zenith distance in Göttingen
by the formula \(\tfrac12(a+a')\) has weight
\(\dfrac{4\alpha\alpha'}{\alpha+\alpha'}\). If the collimation error in Göttingen were known exactly
and were \(=f\),
the determination of the true zenith distance there by the formula
\[
\frac{\alpha(a+f)+\alpha'(a'-f)}{\alpha+\alpha'}
\]

% Printed 33; principal PDF 40; Google PDF 44.
\par\clearpage
\noindent\makebox[12mm][l]{}\makebox[\dimexpr\linewidth-24mm\relax][c]{\small\textbf{III.\quad RESULTS.}}\makebox[12mm][r]{\fontsize{13}{16}\selectfont 33}\par\medskip

\noindent would have weight \(\alpha+\alpha'=\dfrac{4\alpha\alpha'}{\alpha+\alpha'}+
\dfrac{(\alpha-\alpha')^2}{\alpha+\alpha'}\), i.\ e. a greater weight than by the
other method, whenever \(\alpha\) and \(\alpha'\) are unequal. The same holds for
the true zenith distance in Altona, and in this way even observations made
in only one position (such as those of No. 5) would contribute, however slightly,
to improving the precision. The collimation errors at the two sites are
indeed known without \textls[100]{absolute} precision, but it is easy
to see that using the mean values found for them reduces the weight
only very slightly.

\begin{center}\textbf{5.}\end{center}

To obtain a pure result that fully satisfies the requirements of rigorous
theory, however, one must treat the determination of the difference in
latitude, the collimation errors, and the true zenith distances of the
individual stars at one of the two sites as a single problem, in which these unknown
quantities (46 in our case) are to be derived from all the observed quantities
(171) that they determine, by means of as many equations, combined according
to the rules of the calculus of probabilities. If we set the collimation errors
in Göttingen and Altona \(=f\) and \(g\), the difference in latitude \(=h\), and the
true zenith distance of a star in Göttingen \(=k\), the observations of
that star give the four equations with weights
\(\alpha,\alpha',\beta,\beta'\):
\[
\begin{aligned}
a&=k-f\\
a'&=k+f\\
b&=k-g-h\\
b'&=k+g-h.
\end{aligned}
\]

It is scarcely necessary to recall that, to ease the calculation, it is
preferable to introduce, in place of those unknown quantities, the corrections
still to be applied to the values already determined to a close approximation.
If the symbols \(f^0,g^0,h^0,k^0\) denote these approximate values,
one may take
\[
k^0=\frac{\alpha(a+f^0)+\alpha'(a'-f^0)+\beta(b+g^0+h^0)+\beta'(b'-g^0+h^0)}
{\alpha+\alpha'+\beta+\beta'}.
\]

Following that prescription (which should never be lost sight of when applying the method
\par\smallskip
\noindent\textbf{IX.}\hfill 5\hspace{.12\textwidth}\par

% Printed 34; principal PDF 41; Google PDF 45.
\par\clearpage
\noindent\makebox[12mm][l]{\fontsize{13}{16}\selectfont 34}\parbox[t]{\dimexpr\linewidth-24mm\relax}{\centering\scriptsize\bfseries DETERMINATION OF THE DIFFERENCE IN LATITUDE BETWEEN GÖTTINGEN AND ALTONA.}\makebox[12mm][r]{}\par\medskip

\noindent of least squares to cases of even moderate complexity), and using
a suitable \textls[100]{indirect} method of solution, turns a task that would
be intolerably lengthy without the prescription and with direct elimination
into an easy exercise.

\begin{center}\textbf{6.}\end{center}

This calculation, which need not be set out in detail here, yields no
appreciable corrections to the earlier determinations. The correction to
the difference in latitude is \(=-0''{,}014\),
the correction to the collimation error in Göttingen \(=+0''{,}012\), and the
correction to the collimation error in Altona \(=-0''{,}014\). Hence the new
determinations are
\begin{center}
\begin{tabular}{@{}l@{\quad}r@{\enspace}c@{\enspace}l@{}}
Difference in latitude & \(2^\circ0'56''{,}51\), & Weight & \(=217{,}67\)\\
Collimation error in Göttingen & \(3{,}76\) & » & \(\phantom{=}457{,}03\)\\
Collimation error in Altona & \(1{,}39\) & » & \(\phantom{=}437{,}64\).
\end{tabular}
\end{center}

The changes in the true zenith distances of the individual stars in Göttingen
adopted as a basis according to the prescription in the preceding article
are likewise nearly all below \(0''{,}01\). It would be superfluous to list the
resulting values here, since these are the values from which the stellar
declinations given above were derived, on the assumption that the altitude of the
pole at the observing site was \(51^\circ31'47''{,}92\). Here we give, however, the differences
that remain after substituting the values found into the
171 equations.

{\small
\[
\begin{array}{@{}c|c||c|c||c|c||c|c@{}}
\text{Star}&\text{Difference}&\text{Star}&\text{Difference}&
\text{Star}&\text{Difference}&\text{Star}&\text{Difference}\\[.8ex]
1&+0''{,}07&5&+0''{,}03&9&-0''{,}06&13&+0''{,}87\\
&+0{,}09&&-0{,}03&&-0{,}09&&+0{,}30\\
&-0{,}26&&\text{—}&&-0{,}23&&-1{,}35\\
&+0{,}13&&-0{,}03&&+0{,}97&&-0{,}04\\
2&+0{,}56&6&+0{,}62&10&-0{,}71&14&+0{,}40\\
&-0{,}08&&+0{,}03&&-0{,}25&&+0{,}21\\
&-0{,}12&&-1{,}95&&+0{,}70&&-0{,}29\\
&-0{,}53&&-0{,}02&&+0{,}44&&-0{,}48\\
3&+0{,}48&7&-0{,}52&11&+0{,}12&15&-0{,}04\\
&+0{,}34&&+0{,}16&&-1{,}03&&+0{,}19\\
&-0{,}70&&-0{,}08&&+0{,}72&&-0{,}04\\
&-0{,}18&&+0{,}52&&+0{,}71&&-0{,}11\\
4&-0{,}35&8&-0{,}56&12&+0{,}52&16&-0{,}07\\
&-0{,}35&&+0{,}12&&-0{,}87&&-0{,}60\\
&+0{,}79&&+0{,}76&&+0{,}80&&-0{,}05\\
&-0{,}17&&-0{,}23&&-0{,}24&&+0{,}72
\end{array}
\]
\par
}

% Printed 35; principal PDF 42; Google PDF 46.
\par\clearpage
\noindent\makebox[12mm][l]{}\makebox[\dimexpr\linewidth-24mm\relax][c]{\small\textbf{III.\quad RESULTS.}}\makebox[12mm][r]{\fontsize{13}{16}\selectfont 35}\par\medskip

{\small
\[
\begin{array}{@{}c|c||c|c||c|c||c|c@{}}
\text{Star}&\text{Difference}&\text{Star}&\text{Difference}&
\text{Star}&\text{Difference}&\text{Star}&\text{Difference}\\[.8ex]
17&-0''{,}06&24&-0''{,}17&31&-0''{,}41&38&+0''{,}30\\
&-0{,}52&&+0{,}18&&+1{,}11&&-0{,}19\\
&+0{,}96&&+0{,}16&&-0{,}59&&-0{,}24\\
&-0{,}35&&-0{,}20&&+0{,}09&&+0{,}11\\
18&-0{,}23&25&-0{,}03&32&+0{,}14&39&+0{,}90\\
&-0{,}49&&+0{,}22&&-0{,}87&&+1{,}71\\
&+1{,}09&&+0{,}46&&+0{,}63&&-0{,}82\\
&-0{,}34&&-0{,}67&&+0{,}11&&-0{,}69\\
19&+0{,}32&26&+0{,}90&33&+0{,}19&40&-1{,}81\\
&-0{,}36&&+0{,}46&&-0{,}35&&+0{,}75\\
&+0{,}10&&-0{,}06&&+0{,}41&&+0{,}60\\
&+0{,}08&&-1{,}32&&-0{,}16&&-0{,}14\\
20&-0{,}06&27&-0{,}14&34&-0{,}48&41&+0{,}39\\
&-0{,}26&&+1{,}22&&+1{,}17&&-0{,}21\\
&+0{,}66&&-0{,}71&&-0{,}62&&-0{,}38\\
&-0{,}44&&+0{,}46&&-0{,}07&&+0{,}35\\
21&-0{,}22&28&+0{,}46&35&-0{,}08&42&+0{,}09\\
&-0{,}85&&-0{,}53&&-0{,}37&&+0{,}25\\
&+0{,}63&&+0{,}52&&+0{,}25&&-1{,}02\\
&+0{,}37&&-0{,}36&&+0{,}21&&+1{,}11\\
22&+0{,}77&29&-0{,}80&36&-0{,}14&43&-0{,}42\\
&-0{,}55&&-0{,}53&&+1{,}02&&+0{,}86\\
&-1{,}55&&+0{,}58&&-0{,}47&&+0{,}29\\
&+1{,}60&&+0{,}10&&-0{,}58&&-0{,}47\\
23&-0{,}80&30&-0{,}83&37&+0{,}11&&\\
&+0{,}67&&+0{,}74&&+0{,}54&&\\
&-0{,}03&&-0{,}21&&-0{,}93&&\\
&+0{,}03&&+0{,}42&&+0{,}11&&
\end{array}
\]
\par
}

\begin{center}\textbf{7.}\end{center}

The sum of the products of the squares of these 171 differences by the
corresponding numbers of observations is found to be \(=292{,}8249\). By
the theorem already cited (Theoria Comb. Observ. Art. 38), the approximate
value of the mean error of a single observation is the square root of the
fraction whose numerator is that sum and whose denominator is the excess
of the number of observational data compared over the number of unknown
quantities derived from them by the method of least squares, in our case
\(171-46=125\). The resulting mean error is \(=1''{,}5308\), differing little from that
found in Art. 2. The mean error to be expected
in the final result for the
difference in latitude would thus be
\[
=\frac{1''{,}5308}{\surd 217{,}67}=0''{,}1038.
\]
\par\vfill
\noindent\hfill \(5^{*}\)\hspace{.12\textwidth}\par

% Printed 36; principal PDF 43; Google PDF 47.
\par\clearpage
\noindent\makebox[12mm][l]{\fontsize{13}{16}\selectfont 36}\parbox[t]{\dimexpr\linewidth-24mm\relax}{\centering\scriptsize\bfseries DETERMINATION OF THE DIFFERENCE IN LATITUDE BETWEEN GÖTTINGEN AND ALTONA.}\makebox[12mm][r]{}\par\medskip

\begin{center}\textbf{8.}\end{center}

The calculations so far assume that the errors affecting the various
observations can all be regarded as wholly independent of one another
or purely random. This assumption, however, is clearly not quite correct,
since all \(\alpha\) observations used in determining an \(a\) refer, by the
nature of the instrument, to one and the same graduation mark. They therefore
involve, in addition to the purely random errors of observation proper,
the error in the graduation at that mark.
The same holds for \(a'\), \(b\), and \(b'\). The graduation errors are themselves
unknown quantities that may also be regarded as purely random and mutually
independent with respect to the individual 171 observational results, since
the few cases in which several of these referred to the same graduation mark
may be disregarded. Taking this circumstance into account requires a
modification of the above calculations,
although in the end the results
are altogether unaffected for practical purposes.

If the mean error of observation proper, which arises solely from random
causes apart from graduation errors, is denoted by \(m\),
and the mean graduation error by \(\mu\), then the total mean error of an
observation must be set \(=\surd(mm+\mu\mu)\), and the mean error of an
average of \(\alpha\) observations referring to the same graduation mark,
\[
=\ixrad\left(\frac{mm}{\alpha}+\mu\mu\right),
\]
\noindent or, if we set \(\mu\mu=mm\theta\),
\[
=m\ixrad\left(\frac1\alpha+\theta\right).
\]

Thus, insofar as we take the weight of an observation \textls[100]{without graduation error}
as the unit, the weight of \(a\) will now
\[
=\frac{\alpha}{1+\alpha\theta}
\]
\noindent be, and likewise the weights of \(a'\), \(b\), \(b'\), respectively
\[
=\frac{\alpha'}{1+\alpha'\theta},\qquad
\frac{\beta}{1+\beta\theta},\qquad
\frac{\beta'}{1+\beta'\theta}.
\]
With the first method of combination, therefore, the weight of the
\par

% Printed 37; principal PDF 44; Google PDF 48.
\par\clearpage
\noindent\makebox[12mm][l]{}\makebox[\dimexpr\linewidth-24mm\relax][c]{\small\textbf{III.\quad RESULTS.}}\makebox[12mm][r]{\fontsize{13}{16}\selectfont 37}\par\medskip

\noindent result for the difference in latitude from the observations of one star
when we write the previous expression as
\[
\frac4{\dfrac1\alpha+\dfrac1{\alpha'}+\dfrac1\beta+\dfrac1{\beta'}}=p
\]
\noindent must now be taken as
\[
=\frac p{1+p\theta}
\]
\noindent and the mean must be taken from the 42 determinations using these weights.
With the second method of combination, however, it is only necessary
to assign to each of the 171 equations a weight determined
by one of the formulas \(\dfrac{\alpha}{1+\alpha\theta}\), and\ so\ on.

Clearly, a change in the final result itself or in the mean error to be
expected in it can occur only because the new weights are not proportional
to the old ones. The former method merely gave somewhat too much weight
to the results from the longer series of observations. Taking the graduation
errors into account brings their weights closer to equality, the more so the
larger those errors are assumed to be. Thus, in observations with an
instrument whose graduation errors far outweighed
the errors of observation proper, one could only
regard all determinations as equally
reliable.

\begin{center}\textbf{9.}\end{center}

The methods described therefore present no difficulty once the coefficient
\(\theta\) is known. An approximate value for it, which is of course all
that is required here, can be obtained
indirectly.

First, we note that the observations themselves provide a means of determining
the mean error of observation proper, \(m\), with great reliability.
In fact, this error shows up independently of the graduation error in the
differences between the \textls[100]{individual} values of which each \(a\)
(or \(a'\), \(b\), \(b'\)) is the mean, or in their differences from that mean.
If \(\alpha\) were very large, the sum of the squares of the deviations of the
individual values of \(a\) from their mean could be taken as an approx-
\par

% Printed 38; principal PDF 45; Google PDF 49.
\par\clearpage
\noindent\makebox[12mm][l]{\fontsize{13}{16}\selectfont 38}\parbox[t]{\dimexpr\linewidth-24mm\relax}{\centering\scriptsize\bfseries DETERMINATION OF THE DIFFERENCE IN LATITUDE BETWEEN GÖTTINGEN AND ALTONA.}\makebox[12mm][r]{}\par\medskip

\noindent imate determination of \((\alpha-1)mm\). One such determination may
deviate considerably from the correct value in our case, where \(\alpha\) is
never greater than 7. But by the principles of the calculus of probabilities,
the sum of all 171 partial sums (for all \(a\), \(a'\), \(b\), \(b'\) and for all stars)
must differ little from
\[
(\Sigma(\alpha-1)+\Sigma(\alpha'-1)+\Sigma(\beta-1)+\Sigma(\beta'-1))mm,
\]
\noindent which in our case is \(728mm\). We found the sum
of the 171 partial sums to be
\[
=844{,}50
\]
\noindent which gives the very reliable value for \(m\)
\[
=1''{,}0770
\]
\noindent considerably smaller than the value found in Arts. 2 and 7.
The influence of the graduation errors is thus fully confirmed, and accounts
for why the figures obtained previously could not give a pure result.

\begin{center}\textbf{10.}\end{center}

In the absence of direct knowledge of the mean graduation error, \(\theta\)
can now be determined indirectly so that, whether using the first method
with the procedure of Art. 2 or the second method with the procedure of
Art. 7, the mean error of an observation whose weight was taken as the
unit again agrees with the value found for
\(m\).

It did not, however, seem worth repeating such trials until perfect
agreement had been reached. After other considerations had shown that
the final value of \(\theta\) could differ only slightly from \(0{,}2\),
it seemed sufficient to use this value as the basis for the first method
of combination alone, which gave
the following results
\begin{center}
\begin{tabular}{@{}l@{\quad}l@{}}
Difference in latitude & \(=2^\circ0'56''{,}50\)\\
Weight of this determination & \(=104{,}29\)\\
\multicolumn{2}{@{}l@{}}{Mean error of an observation with unit weight,}
\end{tabular}
\end{center}
\[
=1''{,}131
\]

% Printed 39; principal PDF 46; Google PDF 50.
\par\clearpage
\noindent\makebox[12mm][l]{}\makebox[\dimexpr\linewidth-24mm\relax][c]{\small\textbf{III.\quad RESULTS.}}\makebox[12mm][r]{\fontsize{13}{16}\selectfont 39}\par\medskip

\noindent and hence the mean error to be expected in the above final result
\[
=0''{,}1108.
\]

Applying the second method of combination with the same value of
\(\theta\) would probably have given still closer agreement with the above
value of \(m\), perhaps reduced the final result for the difference in latitude
by \(0''{,}01\), and certainly increased the weight of that determination
slightly. It was not considered worth doing the calculation over again on
that account. One may therefore retain the difference in latitude found,
\(2^\circ5'56''50\), and regard its error as probably lying between
the limits \(\pm0''{,}07\).

\begin{center}\textbf{11.}\end{center}

If we retain the above value of \(\theta\), the mean graduation error is
\(=m\surd\theta=0''{,}48\), so the so-called probable graduation error
of the individual marks may be set \(=0''{,}32\). Clearly, this refers
only to the irregular graduation errors, that is, to the deviations of
the individual marks from an imagined uniform graduation fitted to
them as closely as possible. Its
\textls[100]{absolute} accuracy could not properly enter into the question here.
In other words, strictly speaking, the result found for the difference in
latitude, with the precision assigned to it, refers only to mean degrees of
the sector and remains dependent on their absolute accuracy.
The instrument itself gives the astronomer no independent means of
checking \textls[100]{this}. If one considers, however, that the end points of the arc
were laid off by the instrument maker with the greatest care and that only
a small part of the whole arc is involved here, one must grant that this
source can increase the uncertainty of the difference in latitude found
only very slightly. Some check of the absolute accuracy of the graduation
is also afforded by the zenith distances of the same
43 stars which I observed with \textsc{Reichenbach}'s meridian circle.
Their differences from the distances observed with the sector,
when arranged by declination,
show no trace of regularity.

% Printed 40; principal PDF 47; Google PDF 51.
\par\clearpage
\noindent\makebox[12mm][l]{\fontsize{13}{16}\selectfont 40}\parbox[t]{\dimexpr\linewidth-24mm\relax}{\centering\scriptsize\bfseries DETERMINATION OF THE DIFFERENCE IN LATITUDE BETWEEN GÖTTINGEN AND ALTONA.}\makebox[12mm][r]{}\par\medskip

\begin{center}\textbf{12.}\end{center}

The centre of the sector in Göttingen was \(1{,}060\) toises north
and \(7{,}595\) toises east of the centre of the axis of \textsc{Reichenbach}'s
meridian circle. In Altona, by contrast, the centre of the sector was \(13{,}511\)
toises south and \(2{,}578\) west of the centre of the meridian circle there.
The reduction of the difference in latitude between the sector sites to that between
the meridian circles is thus \(0''{,}07\) for Göttingen and \(0''{,}85\) for Altona, and hence
the difference in latitude between the Göttingen and Altona observatories,
referred to the sites of \textsc{Reichenbach}'s meridian circles, is
\[
=2^\circ0'57''{,}42.
\]

\begin{center}\textbf{13.}\end{center}

The absolute altitude of the pole used above to derive the stellar
declinations from the zenith distances rests on 89
observations of Polaris with \textsc{Reichenbach}'s meridian circle, at both
culminations, directly and by reflection from a water surface. Since the
observations from 1824, which make up the greater part, have yet to be
published, I collect all the observations here and remark only that
the direct sighting was usually made as the star reached the
second, fourth (middle), and sixth wires, while the sighting of the
reflected image was made as it reached the first, third, fifth, and
seventh wires. The mean of these zenith distances, reduced to the time of
culmination, is given here with only the refraction removed, using
\textsc{Bessel}'s tables. It therefore still includes the collimation error and the
effect of the flexure of the telescope.

\begin{center}
\textls[100]{Zenith distances of Polaris.}\par\smallskip
1820.\quad Circle east.\par\smallskip
\renewcommand{\arraystretch}{1.12}%
\begin{tabular}{@{}r@{\enspace}l|l@{}}
\ixobs{May 13. Lower culm.}{319^\circ}{50'}{20''{,}73}{220}{5}{3{,}94}{\ixobsCount{3}{Obs.}}{\ixobsCount{4}{»}}\\[.6ex]
\ixobs{13. Upper culm.}{323}{8}{41{,}51}{216}{46}{44{,}31}{\ixobsCount{1}{»}}{\ixobsCount{1}{»}}
\end{tabular}
\end{center}

% Printed 41; principal PDF 48; Google PDF 52.
\par\clearpage
\noindent\makebox[12mm][l]{}\makebox[\dimexpr\linewidth-24mm\relax][c]{\small\textbf{III.\quad RESULTS.}}\makebox[12mm][r]{\fontsize{13}{16}\selectfont 41}\par\medskip

\begin{center}
1824.\quad Circle east.\par\smallskip
\renewcommand{\arraystretch}{1.18}%
\begin{tabular}{@{}r@{\enspace}l|l@{}}
\ixobs{Apr, 20. Upper culm.}{323^\circ}{7'}{52''{,}62}{216}{48}{54{,}93}{\ixobsCount{1}{Obs.}}{\ixobsCount{2}{»}}\\[.7ex]
\ixobs{21. Lower culm.}{319}{52}{30{,}27}{220}{4}{19{,}32}{\ixobsCount{3}{»}}{\ixobsCount{4}{»}}\\[.7ex]
\ixobs{21. Upper culm.}{323}{7}{54{,}16}{216}{48}{54{,}21}{\ixobsCount{3}{»}}{\ixobsCount{4}{»}}\\[.7ex]
\ixobs{25. Lower culm.}{319}{52}{30{,}03}{220}{4}{21{,}10}{\ixobsCount{3}{»}}{\ixobsCount{4}{»}}\\[.7ex]
\ixobs{27. Upper culm.}{323}{7}{55{,}70}{216}{48}{52{,}93}{\ixobsCount{3}{»}}{\ixobsCount{4}{»}}\\[.7ex]
\ixobs{28. Upper culm.}{323}{7}{55{,}40}{216}{48}{52{,}22}{\ixobsCount{3}{»}}{\ixobsCount{4}{»}}\\[.7ex]
\ixobs{29. Lower culm.}{319}{52}{29{,}17}{220}{4}{21{,}34}{\ixobsCount{3}{»}}{\ixobsCount{4}{»}}\\[.7ex]
\ixobs{May 1. Lower culm.}{319}{52}{28{,}59}{220}{4}{22{,}62}{\ixobsCount{3}{»}}{\ixobsCount{4}{»}}\\[.7ex]
\ixobs{1. Upper culm.}{323}{7}{57{,}22}{216}{48}{51{,}66}{\ixobsCount{3}{»}}{\ixobsCount{4}{»}}\\
\multicolumn{2}{@{}c|}{\rule{0pt}{3.3ex}Circle west.}&\\[.5ex]
\ixobs{May 2. Lower culm.}{40^\circ}{4'}{20''{,}00}{139}{52}{27{,}15}{\ixobsCount{3}{Obs.}}{\ixobsCount{4}{»}}\\[.7ex]
\ixobs{8. Upper culm.}{36}{48}{49{,}32}{143}{7}{57{,}63}{\ixobsCount{3}{»}}{\ixobsCount{4}{»}}\\[.7ex]
\ixobs{9. Lower culm.}{40}{4}{22{,}93}{139}{52}{25{,}68}{\ixobsCount{3}{»}}{\ixobsCount{4}{»}}
\end{tabular}
\end{center}

The changes in the declination of Polaris are obtained from \textsc{Bessel}'s
tables as follows:
\par\smallskip
\noindent\textbf{IX.}\hfill 6\hspace{.12\textwidth}\par
\par
\endgroup
